\documentclass[reqno]{amsart}
\usepackage{a4wide}
\usepackage{hyperref}
\usepackage{amssymb}

\AtBeginDocument{{\noindent\small
\emph{Electronic Journal of Differential Equations},
Vol. 2026 (2026), No. 02, pp. 1--15.\newline
ISSN: 1072-6691. URL: https://ejde.math.txstate.edu, DOI: 10.58997/ejde.2026.02}
\thanks{\copyright 2026. This work is licensed under a CC BY 4.0 license.}
\vspace{8mm}}

\begin{document}
\title[\hfilneg EJDE-2026/02\hfil 
Multiplicity of solutions for biharmonic elliptic singular problems]
{Multiplicity of solutions for biharmonic equations with critical
exponent and prescribed singularity}

\author[N. Yagoub, M. E. O. El Mokhtar, A. Matallah, S. Benmansour \hfil EJDE-2026/02\hfilneg]
{Nadjet Yagoub, Mohammed El Mokhtar Ould El Mokhtar, \\
Atika Matallah, Safia Benmansour}

\address{Nadjet Yagoub \newline
Laboratory of Analysis and Control of Partial Differential Equations,
University of Sidi Bel Abbes, Algeria}
\email{nadjet.yagoub@univ-sba.dz}

\address{Mohammed El Mokhtar Ould El Mokhtar \newline
Department of Mathematics,
College of Science,
Qassim University, BO 6644,
Buraidah 51452, Saudi Arabia}
\email{med.mokhtar66@yahoo.fr}

\address{Atika Matallah \newline
Laboratory of Analysis and Control of Partial Differential Equations,
Higher School of Management, Tlemcen, Algeria}
\email{atika\_matallah@yahoo.fr}

\address{Safia Benmansour\newline
Laboratory of Analysis and Control of Partial Differential Equations,
Higher School of Management, Tlemcen, Algeria}
\email{safiabenmansour@hotmail.fr}

\thanks{Submitted September 8, 2025. Published January 8, 2026.}
\subjclass[2020]{35J20, 35IJ60, 47J30}
\keywords{Variational methods; singular potential;  multiplicity of solutions; 
\hfill\break\indent critical Sobolev exponent; biharmonic equation}

\begin{abstract}
In this article, we study the singular critical biharmonic problem
\begin{gather*}
 \Delta^2u-\mu V(x)u=|u|^{2^*-2}u+\lambda f(x) \quad \text{in } \Omega, \\
 u= \frac{\partial u}{\partial n}=0 \quad \text{on } \partial \Omega,
\end{gather*}
where \( \Delta^2\) is the biharmonic operator, \(\Omega\) is an open bounded domain
in \(\mathbb{R}^N\) \((N \geq 5)\) with smooth boundary
\(\partial \Omega, 2^* = \frac{2N}{N-4}\), \(0 < \mu < \bar{\mu} :=
\big( \frac{N(N-4)}{4} \big)^2\), \(f(x)\) and \(V(x)\) are given functions.
By using variational method and Nehari-type constraint, we establish the existence of
multiple solutions for this problem when \(0<\lambda<\lambda^*\), for some
\(\lambda^* > 0\).
\end{abstract}

\maketitle
\numberwithin{equation}{section}
\newtheorem{theorem}{Theorem}[section]
\newtheorem{lemma}[theorem]{Lemma}
\newtheorem{proposition}[theorem]{Proposition}
\newtheorem{remark}[theorem]{Remark}
\allowdisplaybreaks

\section{Introduction}

This article concerns the biharmonic problem
\begin{equation}\label{11}
\begin{gathered}
 \Delta^2u-\mu V(x)u=|u|^{2^*-2}u+\lambda f(x)\quad \text{in }\Omega , \\
 u= \frac{\partial u}{\partial n}=0 \quad \text{on }\partial \Omega ,
 \end{gathered}
\end{equation}
where $\Delta^2$ denotes the biharmonic operator, $\Omega\subset\mathbb{R}^N$
$(N\ge5)$ is a bounded domain with smooth boundary $\partial\Omega$, $\lambda>0$
is a parameter, $2^*=\frac{2N}{N-4}$ is the critical Sobolev exponent and
 $0<\mu<\bar\mu$, where $\bar\mu$ is the best constant for the Rellich inequality
\[
\int_\Omega \frac{u^2}{|x-a|^4}\,dx \leq \frac{1}{\bar\mu}
\int_\Omega |\Delta u|^2\,dx,
\quad \forall a\in\Omega,\; u\in H^2_0(\Omega).
\]
Here, $H^2_0(\Omega)$ denotes the completion of $\mathcal{C}^\infty_0(\Omega)$ with
respect to the norm
\[
\|u\|=\Big(\int_\Omega |\Delta u|^2\,dx\Big)^{1/2}.
\]

In elasticity theory, biharmonic equations with multipolar singular potentials
effectively model thin elastic plates that include \(k\) localized defects or
concentrated loads situated at finitely many points \(a_1, a_2,\dots,a_k\).
These singularities strongly affect deformation and stability and can be used to
eliminate unwanted frequencies or localize vibrational energy
(see Lindsay et al.\ \cite{lin}). Similar models appear in composite and metamaterials, where point inclusions generate stress multipoles influencing effective properties and cloaking effects (see Mao-Huang \cite{met}). In addition, biharmonic singular models can also arise in thin-film physics, surface growth, and nanomechanical resonators (see \cite{ek}, \cite{g}, and \cite{gw}).

In recent years, many authors have studied biharmonic problems, for instance, in the regular case \((\mu=0)\), Qian-Wang in \cite{bh2} proved the existence of at least two distinct solutions for \eqref{11} with \(\mu=0\), \(\lambda=1\) and \(f(x)\) small. In the singular case \((\mu \neq 0)\), we recall that the existence of multiple solutions of \eqref{11} has been studied under different hypotheses on \(V(x)\). For the homogeneous case, D'Ambrossio-Jannelli \cite{j1} considered the Hardy potentiel \(V(x)=|x|^{-4}\) and Kang-Xiong in \cite{b5} established the existence and nonexistence of ground state solutions of \eqref{11} where \(V(x)\) has prescribed singularities of the form \(|x-a_j|^{-4}\) with \(a_j \in \Omega\) and \(j=1, \dots,k\), by using a complicated asymptotic analysis and variational arguments. For the nonhomogenous case, Li et al. \cite{bt} studied \eqref{11} with \(V(x)=|x|^{-s}\), \(\ 0<s\leq 4 \).
Very recently, the authors in \cite{ms} have proved the existence of at least \(2k\)
solutions of the problem
\begin{gather*}
\Delta^2u- \sum_{j=1}^{k} \frac{\mu_j}{|x-a_j|^4}u=|u|^{2^*-2}u
+\sum_{j=1}^{k} \frac{\lambda_j}{|x-a_j|^{4-\alpha_j}}u+ f(x) \quad \text{in }\Omega \\
u= \frac{\partial u}{\partial n}=0 \quad\text{on }\partial \Omega ,
\end{gather*}
where \(\Omega \) is an open bounded domain of
\( \mathbb{R} ^N\) \((N\geq 5)\) with smooth boundary \(\partial \Omega, \) for all
\(j=1, \dots ,k\), \(a_j\in \Omega\) denote the singularity points, \(\lambda_j>0\)
are parameters, \(0<\alpha_j<4\) and \(\mu_j>0\) are real constants satisfying
\(\sum_{j=1}^{k} \mu_j< \bar{\mu}\).
We recall that the problem \eqref{11} with the Laplacian operator was studied by
Chen \cite{c1}. He proved the existence of at least $k$ positive solutions by the
argument developed in \cite{t}. The main goal of this paper is to generalize the
result in \cite{c1} to the biharmonic operator.

From the Rellich inequality, it follows that the best constant
\begin{equation}\label{A}
\mathcal{A}_{\mu}(\Omega)
=\inf_{u\in H^2_0(\Omega)\setminus\{0\}}
\frac{\int_{\Omega}\Big(|\Delta u|^2-\mu\frac{u^2}{|x-a|^4}\Big)\,dx}
{\Big(\int_{\Omega}|u|^{2^*}\,dx\Big)^{2/2^*}},\quad \forall a\in\Omega,\ \mu<\bar\mu,
\end{equation}
is well defined.
Moreover, as shown in \cite{j1,K}, \(\mathcal{A}_\mu(\Omega)\) is independent of the
domain \(\Omega\) and is attained in \(\mathbb{R}^N\) by the family of translated
extremals
\[
\{y_{\varepsilon}(x-a)=\varepsilon^{\frac{4-N}{2}}
U_{\mu}(\varepsilon^{-1}(x-a)), \ \varepsilon>0\},
\]
where \(U_\mu\) is a positive, radially symmetric, and radially decreasing solution
of
\[
\Delta^2 u-\mu\frac{u}{|x|^4}=u^{2^*-1}\quad\text{in }\mathbb{R}^N\setminus\{0\},
\quad u>0.
\]
The normalized translated profiles satisfy
\[
\int_{\mathbb{R}^N}\Big(|\Delta y_{\varepsilon}(x-a)|^{2}
-\mu\frac{|y_{\varepsilon}(x-a)|^{2}}{|x-a|^{4}}\Big)\,dx
=\int_{\mathbb{R}^N}|y_{\varepsilon}(x-a)|^{2^*}\,dx
=\mathcal{A}_{\mu}^{N/4}.
\]
By setting \(\rho=|x|\), the profile \(U_\mu\) has the following sharp
asymptotics:
\begin{gather*}
U_\mu(\rho) = O_1(\rho^{-a(\mu)}) \quad\text{as }\rho \to 0,\\
U_\mu(\rho) = O_1(\rho^{-b(\mu)}), \quad
U_\mu'(\rho) = O_1(\rho^{-b(\mu)-1}) \quad\text{as }\rho \to \infty,
\end{gather*}
where \( \delta=\tfrac{N-4}{2}\) and
\[
a(\mu)=\delta\,\varphi(\mu),\quad b(\mu)=\delta\big(2-\varphi(\mu)\big),
\]
with
\[
\varphi(\mu)=1-\frac{\sqrt{\,N^2-4N+8-4\sqrt{(N-2)^2+\mu}\,}}{N-4},\quad
\mu\in[0,\bar\mu].
\]
In particular, for \(\mu\in[0,\bar\mu)\) one has
\(0\le a(\mu)\le\delta\le b(\mu)\le 2\delta\). There exist positive constants
\(\mathcal{C}_1(\mu),\mathcal{C}_2(\mu)\) such that
\[
0<\mathcal{C}_1(\mu)\le U_\mu(x)\big(|x|^{a(\mu)/\delta}+|x|^{b(\mu)/\delta}\big)^\delta\le\mathcal{C}_2(\mu)
\quad\forall x\in\mathbb{R}^N\setminus\{0\}.
\]

Before stating our main assumptions, we briefly fix some notations.
We denote by \(\mathcal{D}^{2,2}(\mathbb{R}^N)\) the completion of
\(\mathcal{C}_c^{\infty}(\mathbb{R}^N)\) with respect to the norm
\[
 \|u\|_{\mathcal{D}^{2,2}(\mathbb{R}^N)}
 =\Big( \int_{\mathbb{R}^N} |\Delta u|^2 \, dx \Big)^{1/2}.
\]
For \(1 \leq p < \infty\), \(L^p(\Omega)\) denotes the Lebesgue space endowed with
the norm
\[
 \|u\|_{p}=\Big( \int_{\Omega} |u|^p \, dx \Big)^{1/p}.
\]
We write \(\|\cdot\|_{-}\) for the norm in \(H^{-2}(\Omega)\), the dual space of
\(H^2_0(\Omega)\).
The ball of center \(a \in \mathbb{R}^N\) and radius \(r>0\) is denoted by
 \(B(a,r)\). \(l_i\), \(\nu_i\), and \(C_i\) denote positive constants whose values
 are unimportant. For all \(\varepsilon>0, t>0, O(\varepsilon^t)\) denotes the
quantity satisfying \(|O(\varepsilon^t)|/\varepsilon^t\leq C_1\),
\( O_1(\varepsilon^t)\) denotes \(C_2\varepsilon^t \leqslant O_1(\varepsilon^t)
\leqslant C_3\varepsilon^t\) and \(o(\varepsilon^t)\) means
\(o(\varepsilon^t)/\varepsilon^t\to 0\) as
 \(\varepsilon \to 0\) and \(o(1)\) a generic infinitesimal value,
\(\to\) and \(\rightharpoonup\) denote strong and weak convergence,
respectively. Finally, we note
\[ \label{(A0)}
 \mathcal{A}_{0}=\inf\Big\{ \int_{\Omega }\vert \Delta
u\vert ^{2}; \ u\in H^2_0(\Omega), \ \int _{\Omega}|u|^{2^*}=1 \Big\}.
\]

 We now state the structural assumptions used throughout the paper:
\begin{itemize}
\item[(A1)] \( f \in H^{-2}(\Omega) \) and \( f(x) > 0 \) a.e. in \( \Omega \).

\item[(A2)] There exist $k$ different points \( a_1,a_2,\dots,a_k \in \Omega \) such that
\[
 V(x) \in L^{\infty}_{\mathrm{loc}}\big(\Omega \setminus \{a_1,a_2,\dots,a_k\}\big),
 \quad \lim_{x \to a_j} V(x)\,|x-a_j|^4 = 1.
\]
Moreover, there exist \( \delta_0 > 0 \) and \( \alpha, \beta > 2\big(b(\mu)
- \delta\big) > 0 \) such that, for all \( x \in B(a_j, \delta_0) \) and
\( j \in \{1,\dots,k\} \),
\begin{equation}\label{v}
 1 - |x-a_j|^{\beta} \ \leq\ |x-a_j|^4 V(x) \ \leq\ 1 - |x-a_j|^{\alpha}.
\end{equation}
Here, \( \delta_0 \) is chosen so that \( |a_i-a_j| \geq 4\delta_0 \) for
\( i \neq j \) and \( B(a_j, \delta_0) \subset \Omega \).

\item[(A3)] There exists a constant \( 0 < C < 1 \) such that
\[
 \mu \int_{\Omega} V(x) u^2\,dx \ \leq\ C \int_{\Omega} |\Delta u|^2\,dx,
 \quad \forall u \in H^2_0(\Omega).
\]
\end{itemize}
We are ready to state our main result.

 \begin{theorem}\label{tm}
Assume that {\rm (A1)--(A3)} are satisfied and \(0<\mu<\bar{\mu}\).
Then there exists \(\lambda ^*> 0\) such that for all \(0<\lambda <\lambda^*\),
problem \eqref{11} has at least \(k\) solutions on \(H^2_0(\Omega)\).
\end{theorem}

\begin{remark} \rm
It is worth noting that, unlike \cite{ms}, which studied nonhomogeneous biharmonic
problem with Rellich-type singularities \(V(x)= \sum_{j=1}^{k} |x-a_j|^{-4}\) with
a classical perturbation \(\sum_{j=1}^{k} |x-a_j|^{\alpha_j-4}\), our work considers
a more rigid class of multiple singular potentials \(V(x)\).
Specifically, we impose sharper two-sided bounds near each pole and uniform separation
between the singularities, making the problem \eqref{11} more interesting and delicate.
\end{remark}

This article is organized as follows. In Section 2, we give some preliminaries.
In Section 3, we present the proofs of several technical lemmas and propositions.
Section 4 is devoted to the proof of Theorem \ref{tm}.

\section{Preliminary results}

To prove Theorem \ref{tm}, we will use critical point theory.
On \(H^2_0(\Omega)\), we define the energy functional associated with the problem
\eqref{11} by
\[
 J_\mu (u)=\frac{1}{2}\int_{\Omega}(|\Delta u|^2-\mu V(x)u^2)\,dx
 -\frac{1}{2^*} \int_{\Omega} |u|^{2^*}\,dx-\lambda \int_{\Omega} fu\,dx.
\]
We say that \(u\) is a weak solution of \eqref{11} if \(u\in H^2_0(\Omega)\) and for
all \( \varphi \in H^2_0(\Omega)\), we have
\[
\langle J'_\mu (u), \varphi\rangle=\int_{\Omega}
 (\Delta u\Delta \varphi-\mu V(x)u\varphi)\,dx- \int_{\Omega} |u|^{2^*-2}u\varphi\,dx
-\lambda \int_{\Omega} f \varphi\,dx=0.
\]
We define
\[
 \mathcal{N}_\mu=\{ u \in H^2_0(\Omega): u\neq 0, \ \langle J'_\mu (u),
 u\rangle=0 \}.
\]
First, we give some energy estimates.

\begin{lemma}\label{lm}
 Let $u$ be a solution of \eqref{11}, then for any $\lambda>0$,
$$
 J_{\mu}(u)\geq -\chi \lambda^2,
$$
where
$$
\chi =\frac{(N+4)^2}{32N(1-C)}\|f\|_-^2.
$$
 \end{lemma}

\begin{proof}
 Let \(u\) be a solution of \eqref{11}, then we have
\begin{align*}
J_\mu (u)
&=\frac{1}{2}\int_{\Omega}(|\Delta u|^2-\mu V(x)u^2)\,dx-\frac{1}{2^*} \int_{\Omega} |u|^{2^*}\,dx-\lambda \int_{\Omega} fu\,dx\\
&=\frac{2}{N}\int_{\Omega}(|\Delta u|^2-\mu V(x)u^2)\,dx-\lambda\frac{N+4}{2N} \int_{\Omega} fu\,dx\\
&\geq \frac{2}{N}(1-C)\|u\|^2-\lambda \frac{N+4}{2N}\|u\|\|f\|_-\,.
\end{align*}
 For \(t>0\), we set
 \[
 \bar{h}(t)= \frac{2}{N}(1-C)t^2-\lambda \frac{N+4}{2N}\|f\|_-t,
 \]
 then, we obtain
\[
 \bar{h}(t)\geq \bar{h}(\bar{t})=-\lambda^2\frac{(N+4)^2}{32N(1-C)}\|f\|_-,
\]
where
 \(\bar {t}=\lambda \frac{N+4}{8(1-C)}\|f\|_-\). This completes the proof.
\end{proof}

 Next, we define \(\psi _\mu: (0, +\infty)\to \mathbb{R}\) by 
 \(\psi _\mu(t)=\langle J_{\mu}'(tu), tu\rangle\), that is
\[
 \psi _\mu(t)= t^2\int_{\Omega}(|\Delta u|^2-\mu V(x)u^2)\,dx-t^{2^*} \int_{\Omega}|u|^{2^*}\,dx-\lambda t \int_{\Omega}fu\,dx,
\]
for all \(u\in \mathcal{N}_\mu\). So
\begin{align*}
 \psi' _\mu(1) 
 &=2 \int_{\Omega}\left( |\Delta u|^2-\mu V(x)u^2\right)\,dx - 2^*\int_{\Omega}|u|^{2^*}\,dx -\lambda \int_{\Omega}fu\,dx, \\ \label{1}
 &=\int_{\Omega}\left( |\Delta u|^2-\mu V(x)u^2\right)\,dx -(2^*-1)\int_{\Omega}|u|^{2^*}\,dx.
\end{align*}
We split \(\mathcal{N}_\mu\) into three parts
\begin{gather*}
 \mathcal{N}_\mu^+= \{u\in N_\mu: \psi' _\mu(1) >0\},\\
\mathcal{N}_\mu^0= \{u\in N_\mu: \psi' _\mu(1) =0\}, \\
 \mathcal{N}_\mu^-= \{u\in N_\mu: \psi' _\mu(1)<0\}. 
\end{gather*}
We now derive some basic properties of \(\mathcal{N}_\mu^+, \mathcal{N}_\mu^0\) and 
\(\mathcal{N}_\mu^-\).

\begin{lemma}\label{l2}
Assume that {\rm (A1)--(A3)} are satisfied and \(0<\mu<\bar{\mu}\). 
Then there exists \(\lambda_1>0\) such that for any \(\lambda \in (0,\lambda_1)\), 
\(\mathcal{N}_\mu^0=\emptyset \ \ and \ \ \mathcal{N}_\mu^- \neq\emptyset\).
\end{lemma}

 \begin{proof}
Arguing by contradiction, we assume that there are \(\lambda_n\to 0\) such that 
\(\mathcal{N}_\mu^0\neq \emptyset\), then
\begin{equation}\label{s1}
 \int_{\Omega}(|\Delta u_n|^2-\mu V(x)u_n^2)\,dx=(2^*-1) \int_\Omega |u_n|^{2^*}\,dx,
\end{equation}
and
\begin{equation}\label{s2}
\int_{\Omega}(|\Delta u_n|^2-\mu V(x)u_n^2)\,dx= \int_\Omega |u_n|^{2^*}\,dx+\lambda_n \int_\Omega f u_n\,dx,
\end{equation}
by Sobolev inequality, (A3) and \eqref{s1}, we obtain
\[
 \mathcal{A}_0^{-2^*/2}\Big( \int_{\Omega}|\Delta u_n|^2\,dx\Big)^{2^*/2} 
 \geq \int_{\Omega}|u_n|^{2^*}\,dx
 \geq \frac{1-C}{2^*-1} \int_{\Omega}|\Delta u_n|^2\,dx,
\]
and so
\begin{equation}\label{2}
 \int_{\Omega}|\Delta u_n|^2\,dx\geq \tilde{C} \mathcal{A}_0^{N/4},
\end{equation}
with \(\tilde{C}=\big( \frac{1-C}{2^*-1}\big) ^{\frac{N-4}{4}}\).
Combining \eqref{s1}, \eqref{s2} with \eqref{2}, we obtain
\begin{align*}
 0&=\Big( 1-\frac{1}{2^*-1}\Big) \int_{\Omega}(|\Delta u_n|^2-\mu V(x)u_n^2)\,dx-\lambda_n\int_{\Omega}fu_n\,dx\\
 &\geq \Big( 1-\frac{1}{2^*-1}\Big)(1-C) \int_{\Omega}|\Delta u_n|^2\,dx-\lambda_{n}\|u_n\| \|f\|_->0.
 \end{align*}
Since \(\lambda_n\to 0\). This contradiction implies that there is \(\bar{\lambda}>0\) 
such that \(\mathcal{N}_\mu^0=\emptyset\) for any \(\lambda\in (0,\bar{\lambda})\). \\
 Now, for \(u\in H^2_0(\Omega)\setminus\{0\}, t>0\), we consider the function
 \[
 g(t)= t\int_{\Omega}(|\Delta u|^2-\mu V(x)u^2)\,dx
 -t^{2^*-1} \int_{\Omega}|u|^{2^*}\,dx-\lambda \int_{\Omega}fu\,dx.
 \]
 Let
 \[
 t_{\rm max}=\Big( \frac{\int_{\Omega}(|\Delta u|^2-\mu V(x)u^2)\,dx}
 {(2^*-1)\int_{\Omega}|u|^{2^*}\,dx}\Big)^{ \frac{1}{2^*-2}}.
 \]
 It is clear that \(g(t)\) achieves its maximum at \(t_{\rm max}\) and we have
 \begin{align*}
 g(t_{\rm max})
 &=\Big( \Big( \frac{1}{2^*-1}\Big) ^{\frac{1}{2^*-2}}
  -\Big( \frac{1}{2^*-1}\Big) ^{\frac{2^*-1}{2^*-2}}\Big) 
  \frac{\big(\int_{\Omega}(|\Delta u|^2-\mu V(x)u^2)\,dx\big)^{\frac{2^*-1}{2^*-2}} }
  {\|u\|_{2^*}^{\frac{2^*}{2^*-2}}}
  -\lambda\int_{\Omega}fu\,dx\\
 &=C_N\frac{\left( \int_{\Omega}(|\Delta u|^2-\mu V(x)u^2)\,dx
  \right) ^{\frac{N+4}{8}}}{\|u\|_{2^*}^{N/4}}-\lambda\int_{\Omega}fu\,dx\\
 &\geq C_N \frac{\left( (1-C) \int_{\Omega}|\Delta u|^2\,dx\right) ^{\frac{N+4}{8}}}{\left( \frac{\|u\|}{\sqrt{\mathcal{A}_0}}\right) ^{N/4}}-\lambda\int_{\Omega}fu\,dx\\
 &\geq C_N (1-C)^{\frac{N+4}{8}} \|u\|( \sqrt{\mathcal{A}_0}) ^{N/4}
  -\lambda \|u\|\|f\|_-\,,
 \end{align*}
where \(C_N= \big( \frac{1}{2^*-1}\big) ^{\frac{1}{2^*-2}}-\left( \frac{1}{2^*-1}\right) ^{\frac{2^*-1}{2^*-2}}\).\\
 Let
 \[
 \bar {\bar {\lambda}}=\frac{ C_N (1-C)^{\frac{N+4}{8}}
 \big( \sqrt{\mathcal{A}_0}\big) ^{N/4}}{\|f\|_-}>0,
 \]
then for \(\lambda\in (0,\bar{\bar {\lambda}})\), one has that \(g(t_{\rm max})>0\).
Then there exists \(t^+=t^+(u)\) such that \(0<t_{\rm max}<t^+, g(t^+)=0\) and 
\(g'(t^+)<0\), it follows that \(t^+u\in \mathcal{N}_\mu^-\). In conclusion, for \(\lambda_1=\min \{\bar {\lambda}, \bar {\bar {\lambda}}\}\) we have
 \begin{equation*}
 \mathcal{N}_\mu^0=\emptyset \quad\text{and}\quad
 \mathcal{N}_\mu^-\neq\emptyset. \qedhere
 \end{equation*}
 \end{proof}

 \section{Localization of constraints}
 
We minimize the functional \(J_\mu\) on some subsets of
 constraints \(\mathcal{N}_\mu\). For this similar to \cite{bh}, we define a map of 
 ``Barycenter type'' \(\beta_j: H^2_0(\Omega)\setminus\{0\} \to \mathbb{R}^N\), as
\[
 \beta_{j}(u)=\frac{\int_{\Omega} \psi_{j}(x)|\Delta u|^2\,dx}
 {\int_{\Omega} |\Delta u|^2\,dx}, \quad \text{for }  j\in \{1,2,\dots,k\}
\]
where \(\psi_{j}(x)= \min \{\delta_0, |x-a_j|\}\). For \(r_0=\frac{\delta_0}{3}\),
we set
\[
 \mathcal{N}_j^-=\{ u\in \mathcal{N}_\mu^-, \beta_j(u)<r_0 \} ,\quad
 \Gamma_j^-=\{ u\in \mathcal{N}_\mu^-, \beta_j(u)=r_0 \}.
\]
From \cite{c1} and \cite{hw}, we derive the following result.
 
\begin{proposition}\label{p1}
Let \(r_0\) be as above, if \(\beta_j(u)<r_0\) then
\[
 \int_{\Omega} |\Delta u|^2\,dx \geq 3 \int_{\Omega \setminus B(a_j, r_0)} |\Delta u|^2\,dx.
\]
\end{proposition}

 \begin{remark}\label{RM} \rm
We deduce from Proposition \ref{p1}, that if 
\(u\in H^2_0(\Omega)\setminus\{0\}, \beta_j(u)\leq r_0\)  and \(\beta_i(u)\leq r_0\),
then \(i=j\). Indeed, notice that if \( \beta_j(u)\leq r_0, \ \beta_i(u)\leq r_0\) 
and \(j\neq i\), then by Proposition \ref{p1}, we have
 \[
 2\int_{\Omega} |\Delta u|^2\,dx 
 \geq 3 \Big( \int_{\Omega \setminus B(a_j, r_0)} |\Delta u|^2\,dx
 +\int_{\Omega \setminus B(a_i, r_0)} |\Delta u|^2\,dx\Big) 
 \geq 3\int_{\Omega} |\Delta u|^2\,dx,
 \]
which is a contradiction.
\end{remark}

 Now, for \(j\in \{1,2,\dots,k\}\), we consider the following variational problems
\[
 c_{\lambda, j}=\inf \{J_{\mu}(u), u\in \mathcal{N}_j^-\}, \quad
  \hat{c}_{\lambda, j}=\inf \{J_{\mu}(u), u\in \Gamma_j^-\}.
\]
For \(j\in \{1,2,\dots,k\}\) fixed, choose the radial cut-off function 
\(\phi(x)=\phi(|x|)\in \mathcal{C}^{\infty}_0(B(0,2\delta _0))\) such that 
\(0\leq \phi(x)\leq 1\) in \(B(0, 2\delta_0)\) and \(\phi(x)=1\) in \(B(0, \delta_0)\).

We set \(u_{\varepsilon, j}(x)= \phi(x-a_j) y_\varepsilon(x-a_j)\). 
The following asymptotic properties hold.

\begin{lemma}
Assume that \(N\geq 5\) and \(0<\mu<\bar {\mu}\). Then, as \(\varepsilon \to 0\), we
have the following estimates
\begin{gather} \label{a0}
 \int_{\Omega}\Big( |\Delta u_{\varepsilon, j}|^2-\mu \frac{|u_{\varepsilon, j}
  |^2}{|x-a_j|^4}\Big)\,dx =\mathcal{A}_{\mu}^{N/4}+O(\varepsilon^{2(b(\mu)
  -\delta)}), \\ 
  \label{a1}
  \int_{\Omega}| u_{\varepsilon, j}|^{2^*}\,dx =\mathcal{A}_{\mu}^{N/4}
  +O(\varepsilon^{2^*(b(\mu)-\delta)}), \\ 
  \label{a2}
  \int_{\Omega} |x-a_j|^{\alpha-4}| u_{\varepsilon, j}|^2\,dx 
  =O(\varepsilon^{2(b(\mu)-\delta)}), \quad \text{since }  \alpha> 2(b(\mu)-\delta).
\end{gather}
\end{lemma}

 \begin{proof}
 The proof is similar to \cite{j1} and \cite{K}.
 \end{proof}
 \begin{proposition}\label{PP}
 Let $(A1), (A2)$ and $(A3)$ be verified, and \(0<\mu<\bar{\mu}\). Then, there exists \(\lambda_2>0\), such that for all \(\lambda\in(0,\lambda_2)\) there holds
\begin{equation}\label{s3}
 c_{\lambda, j}<\frac{2}{N}\mathcal{A}_{\mu}^{N/4}-\chi \lambda^2.
\end{equation}
 \end{proposition}
 \begin{proof}
 As \(u_{\varepsilon, j}(x)\neq 0\), from Lemma \ref{l2} we can find 
 \(t_{\varepsilon, j}^-=t^-(u_{\varepsilon, j})>0\), such that for any 
 \(\lambda \in (0,\lambda_1)\) we have
  \(t_{\varepsilon, j}^-u_{\varepsilon, j}\in \mathcal{N}^-_j\). Since
 \[
 \beta_{j}(t_{\varepsilon, j}^-u_{\varepsilon, j})
 =\frac{\int_{\Omega}\psi_j(x)|\Delta u_{\varepsilon, j}|^2\,dx}
 {\int_{\Omega}|\Delta u_{\varepsilon, j}|^2\,dx}\to
 0, \quad\text{as}\quad \varepsilon\to 0,
 \]
it follows that there exists \(\varepsilon_1>0\) such that 
\(\beta_{j}(t_{\varepsilon, j}^-u_{\varepsilon, j})<r_0\) for any 
\(\varepsilon \in(0,\varepsilon_1)\), that is 
\(t_{\varepsilon, j}^-u_{\varepsilon, j}\in \mathcal{N}^-_j\). Hence, we have
\[
 \sup_{t> t_{\rm max}}J_{\mu}(tu_{\varepsilon, j})\leq \sup_{t>0}
 \Big\{ \frac{t^2}{2}\int_{\Omega}\left( |\Delta u_{\varepsilon, j}
 |^2-\mu V(x)u_{\varepsilon, j}^2\right)\,dx 
 -\frac{t^{2^*}}{2^*}\int_{\Omega}|u_{\varepsilon, j}|^{2^*}\,dx \Big\} 
  -t_{\rm max}\lambda \int_{\Omega}fu_{\varepsilon, j}\,dx,
\]
 for $t>0$, we consider the function
 \[
 \bar {g}(t)=\frac{t^2}{2}\int_{\Omega}\left( |\Delta u_{\varepsilon, j}|^2
 -\mu V(x)u_{\varepsilon, j}^2\right)\,dx
 -\frac{t^{2^*}}{2^*}\int_{\Omega}|u_{\varepsilon, j}|^{2^*}\,dx.
 \]
 Using \eqref{v}, we obtain
\[
 \bar {g}(t)\leq \frac{t^2}{2}\Big[ \int_{\Omega}\Big( |\Delta u_{\varepsilon, j}|^2
 -\mu \frac{u_{\varepsilon, j}^2}{|x-a_j|^4}\Big)\,dx
 +\mu \int_{\Omega}|x-a_j|^{\beta-4}u_{\varepsilon, j}^2\,dx\Big] 
  -\frac{t^{2^*}}{2^*}\int_{\Omega}|u_{\varepsilon, j}|^{2^*}\,dx.
\]
 On the other hand, from
 \[
 \sup_{t\geq 0}\Big( \frac{t^2}{2}B_1-\frac{t^{2^*}}{2^*}B_2\Big)
  =\frac{2}{N}B_1^{N/4}B_2^{\frac{4-N}{4}}; \quad  B_1, B_2>0,
 \]
 and  \eqref{a0}, \eqref{a1} and \eqref{a2}, we obtain
 \[
 \sup_{t\geq 0}\bar {g}(t)\leq \frac{2}{N}\mathcal{A}_{\mu}^{N/4}+o(\varepsilon^{2(b(\mu)-\delta)}).
 \]
Hence, for all \(0<\varepsilon<\varepsilon_1\), we have
\begin{align*}
 \sup_{t> t_{\rm max}}J_{\mu}(tu_{\varepsilon, j})
 &\leq \frac{2}{N}\mathcal{A}_{\mu}^{N/4}+o(\varepsilon^{2(b(\mu)-\delta)})-t_{\rm max}\lambda \int_{\Omega}fu_{\varepsilon, j}\,dx\\
 &\leq \frac{2}{N}\mathcal{A}_{\mu}^{N/4}+o(\varepsilon^{2(b(\mu)-\delta)})-C_4\lambda \varepsilon^{b(\mu)-\delta}
\end{align*}
for some positive constant $C_4$ independent of $\varepsilon$ and $j$. We write
\[
 o(\varepsilon^{2(b(\mu)-\delta)})=k(\varepsilon )\varepsilon^{2(b(\mu)-\delta)},
\]
where \(k(\varepsilon )\to 0\) as \(\varepsilon \to 0\). Taking
\begin{equation}\label{s4}
 \varepsilon ^{2(b(\mu)-\delta)}=(\frac{2\chi}{C_1})^2\lambda^2, \quad
  k(\varepsilon)< \chi(\frac{C_1}{2\chi})^2,
\end{equation}
and choosing \(\lambda_2\) small enough, we have taht
 \(\varepsilon< \varepsilon_1\) 
for all \(\lambda\in (0,\lambda_2)\). 
Substituting \(\varepsilon\) as above to \eqref{s4} gives \eqref{s3}.
 \end{proof}

\begin{proposition}\label{p3}
For \(j\in \{1,\dots,k\}\). Let {\rm (A1)--(A3)} be satisfied and 
\(0<\mu<\bar{\mu}\). Then there exists \(\lambda_3>0\), such that for all
\(\lambda\in(0,\lambda_3)\)
\[
 \hat{c}_{\lambda, j}> \frac{2}{N}\mathcal{A}_{\mu}^{N/4}.
\]
\end{proposition}
 
\begin{proof}
We argue by contradiction. Suppose that there exists a sequence \(\lambda_n\to 0\) 
such that
 \[
 \hat{c}_{\lambda_n,j}\to c\leq \frac{2}{N}\mathcal{A}_{\mu}^{N/4}
 \]
for some \(1\leq j\leq k\). Hence, we can find a sequence 
\(\{u_n\}_n\subset H^2_0(\Omega)\) with \(u_n\in \mathcal{N}^-_j\) satisfying
\begin{equation}\label{d1}
 J_\mu(u_n)\to c\leq \frac{2}{N}\mathcal{A}_{\mu}^{N/4}, \quad\text{as }
  n\to \infty.
\end{equation}
From $\{u_n\}_n\subset \mathcal{N}^-_j$, we have
\begin{equation}\label{d2}
 \int_{\Omega}(|\Delta u_n|^2-\mu V(x)u_n^2)\,dx
 - \int_{\Omega} |u_n|^{2^*}\,dx=\lambda_n \int_{\Omega} fu_n \,dx=o(1).
 \end{equation}
This means that
\begin{align*}
 c+o(1)&=J_\mu(u_n)\\
 &=\frac{1}{2}\int_{\Omega}(|\Delta u_n|^2-\mu V(x)u_n^2)\,dx-\frac{1}{2^*} \int_{\Omega} |u_n|^{2^*}\,dx+o(1)\\
 &=\frac{2}{N}\int_{\Omega}(|\Delta u_n|^2-\mu V(x)u_n^2)\,dx+o(1),
\end{align*}
and
 \begin{align*}
 \int_{\Omega}(|\Delta u_n|^2-\mu V(x)u_n^2)\,dx&=\int_{\Omega} |u_n|^{2^*}\,dx+o(1)\\
 &\leq \Big( \frac{\int_{\Omega}|\Delta u_n|^2\,dx}{\mathcal{A}_0}\Big) ^{2*/2}+o(1)\\
 &\leq \Big( \frac{\int_{\Omega}(|\Delta u_n|^2-\mu V(x)u_n^2)\,dx}{(1-C)\mathcal{A}_0}
 \Big) ^{2^*/2}+o(1),
 \end{align*}
 which together imply that
 \[
 Nc+1\geq \int_{\Omega}(|\Delta u_n|^2-\mu V(x)u_n^2)\,dx\geq \left( (1-C)\mathcal{A}_0\right) ^{N/4}>0
 \]
holds for \(n\) large. From (A3), we obtain
\[
 (1-C)\|u_n\|^2\leq \int_{\Omega}(|\Delta u_n|^2-\mu V(x)u_n^2)\,dx \leq C'\|u_n\|^2,
\]
which provides that
\begin{equation}\label{d3}
 0<\nu_1\leq \|u_n\|\leq \nu_2
\end{equation}
for some positive constants $\nu_1$ and $\nu_2$. From \eqref{d2}, we set
\[
 \lim_{n\to \infty}\int_{\Omega}(|\Delta u_n|^2-\mu V(x)u_n^2)\,dx
 =\lim_{n\to \infty} \int_{\Omega} |u_n|^{2^*}\,dx=l_1>0.
\]
We define the Levy concentration functions
\[
 \varrho_n(r)=\sup_{y\in \Omega}\int _{B(y, r)}|u_n|^{2^*}\,dx,
\]
since, for every $n$ we have $r_n>0$ such that \(\varrho_n(r_n)=\frac{l_1}{2}\), then there exists \(y_n\in \Omega\) such that
 \[
 \int _{B(y_n, r_n)}|u_n|^{2^*}\,dx=\varrho_n(r_n)=\frac{l_1}{2}.
 \]
 Let us define the rescaled functions by
\[
 w_n(\tilde{x})=r_n^{\frac{N-4}{2}}u_n(r_n\tilde{x}+y_n),
\]
then \(w_n(\tilde{x})\in H^2_0(\Omega_n)\) with \(\Omega_n=(\Omega-y_n)/r_n\). By extending \(w_n(\tilde{x})\) to be zero for \(\tilde{x}\) outside \(\Omega_n\), we obtain \(w_n(\tilde{x})\in \mathcal{D}^{2,2}(\mathbb{R}^N)\), with
\begin{align*}
 \int_{\mathbb{R}^N}|\Delta w_n(\tilde{x})|^2\,d\tilde{x}&=\int_{\Omega_{n}}|r_n^{\frac{N-4}{2}}\Delta u_n(r_n\tilde{x}+y_n)|^2\,d\tilde{x}\\
 &=\int_{\Omega_{n}}r_n^N|\Delta u_n(r_n\tilde{x}+y_n)|^2\,d\tilde{x}\\
 &=\int_{\Omega}|\Delta u_n(x)|^2\,dx.
\end{align*}
Similarly, we obtain
\begin{gather*}
\int_{\mathbb{R}^N}r_n^4V(r_n\tilde{x}+y_n)w_n^2(\tilde{x})\,d\tilde{x}
=\int_{\Omega}V(x)u_n^2(x)\,dx, \\
 \int _{\mathbb{R}^N}|w_n(\tilde{x})|^{2^*}\,d\tilde{x}
 = \int_{\Omega}|u_n(x)|^{2^*}\,dx.
\end{gather*}
Then by \eqref{d3}, \(\{w_n\}_n\) is bounded in \(\mathcal{D}^{2,2}(\mathbb{R}^N)\),
thus we can assume that
 \[
 w_n(\tilde{x})\rightharpoonup w_0(\tilde{x}) \text{ in }
  \mathcal{D}^{2,2}(\mathbb{R}^N), \quad
 w_n(\tilde{x})\to w_0(\tilde{x}) \text{ a.e  in } \mathbb{R}^N.
 \]
Set \(\tilde{w}_n(\tilde{x})= w_n(\tilde{x})-w_0(\tilde{x})\), we have
from \eqref{d2}
\begin{equation}\label{d4}
 \int_{\mathbb{R}^N}\left( |\Delta w_n(\tilde{x})|^2
 -r_n^4 V(r_n\tilde{x}+y_n)w_n^2(\tilde{x})\right)\, d\tilde{x}-\int _{\mathbb{R}^N}|w_n(\tilde{x})|^{2^*}\,d\tilde{x}=o(1),
\end{equation}
by using Brezis-Lieb lemma, we obtain
 \begin{gather*}
 \int_{\mathbb{R}^N}|\Delta \tilde{w}(\tilde{x})|^2\,d\tilde{x}=\int_{\mathbb{R}^N}|\Delta w_n(\tilde{x})|^2\,d\tilde{x}-\int_{\mathbb{R}^N}|\Delta w_0(\tilde{x})|^2\,d\tilde{x}+o(1)\\
 \int _{\mathbb{R}^N}|\tilde{w}_n(\tilde{x})|^{2^*}\,d\tilde{x}=\int _{\mathbb{R}^N}|w_n(\tilde{x})|^{2^*}\,d\tilde{x}-\int _{\mathbb{R}^N}|w_0(\tilde{x})|^{2^*}\,d\tilde{x}+o(1)\\
 \int_{\mathbb{R}^N}r_n^4V(r_n\tilde{x}+y_n)\tilde{w}_n^2(\tilde{x})\,d\tilde{x}
 =\int_{\mathbb{R}^N}r_n^4V(r_n\tilde{x}+y_n)w_n^2(\tilde{x})\,d\tilde{x}
  -\int_{\mathbb{R}^N}r_n^4V(r_n\tilde{x}+y_n)w_0^2(\tilde{x})\,d\tilde{x}+o(1).
 \end{gather*}
Recalling that \(r_n\tilde{x}+y_n\in \Omega\) and \(\Omega\) is bounded, we can assume
up to a subsequence, that \(r_n\to \bar{r}\geq 0\) and
\(y_n \to \bar {y}\in \overline{\Omega}\), so we will distinguish two cases. Before
doing so, we need to clarify one notation.
Whenever writing \( r_n\tilde{x}+y_n \in (or \notin) B(a_j, \delta_0), \) always
mean that there is a natural number \(N_1\) such that for all \(n>N_1\),
there holds \( r_n\tilde{x}+y_n \in (or \notin) B(a_j, \delta_0), \) for any given
\(\tilde{x}\).
\smallskip

\noindent\textbf{Case (I):} 
If \( r_n\tilde{x}+y_n \notin B(a_j, \delta_0)\), then by taking into account the 
definition of \(\psi_j(x)\), we obtain
\[
 r_0=\beta_j(u_n)=\frac{\int_{\Omega} \psi_{j}(x)|\Delta u_n(x)|^2\,dx}
 {\int_{\Omega} |\Delta u_n(x)|^2\,dx}
 =\frac{\int_{\mathbb{R}^N} \psi_{j}(r_n\tilde{x}+y_n)|\Delta w_n(\tilde{x})|^2
 \,d\tilde{x}}{\int_{\mathbb{R}^N} |\Delta w_n(\tilde{x})|^2\,d\tilde{x}} \to \delta_0,
\]
which is a contradiction to the choice of \(r_0\).
\smallskip

\noindent \textbf{Case (II):} When \(r_n\tilde{x}+y_n \in B(a_j, \delta_0)\), we 
distinguish three steps.
\smallskip

\noindent\textbf{Step 1.} If \(r_n\to \bar{r}=0\) and \(\bar {y}=a_j\), then we have
 \[
 r_0=\beta_j(u_n)=\frac{\int_{\mathbb{R}^N} \psi_{j}(r_n\tilde{x}+y_n)|\Delta
  w_n(\tilde{x})|^2\,d\tilde{x}}
 {\int_{\mathbb{R}^N} |\Delta w_n(\tilde{x})|^2\,d\tilde{x}} \to \psi_j(a_j)=0, 
 \quad\text{as } n\to \infty,
 \]
which is a contradiction.
\smallskip

\noindent\textbf{Step 2.} If \(r_n\to \bar{r}>0\), we consider the following two 
sub-steps:
\smallskip

 \textbf{Sub-step 2.1.}
  If \(\|\tilde{w}_n\|_{ \mathcal{D}^{2,2}(\mathbb{R}^N)}\to 0\), from \eqref{A} and 
 \eqref{d4}, we have
 \begin{align*}
&\int_{\mathbb{R}^N}\Big( |\Delta w_n(\tilde{x})|^2-\mu \frac{r_n^4}{\left| 
r_{n}\tilde{x}+y_n-a_j\right| ^4}w_n^2(\tilde{x})\Big)\,d\tilde{x}\\
 &\geq \mathcal{A}_{\mu} \Big( \int _{\mathbb{R}^N}|w_n(\tilde{x})|^{2^*}
  \,d\tilde{x}\Big) ^{2/2^*}\\
 &=\mathcal{A}_{\mu} \Big(\int_{\mathbb{R}^N}\left( |\Delta w_n(\tilde{x})|^2
  -\mu r_n^4 V(r_n\tilde{x}+y_n)w_n^2(\tilde{x})\right) \,d\tilde{x}\Big)^{2/2^*}\\
 &\geq \mathcal{A}_{\mu} \Big(\int_{\mathbb{R}^N}\Big( |\Delta w_n(\tilde{x})|^2
 -\mu \frac{r_n^4}{\left| r_{n}\tilde{x}+y_n-a_j\right| ^4}w_n^2(\tilde{x})\Big) 
 \,d\tilde{x}\Big)^{2/2^*}.
 \end{align*}
Hence, we obtain
 \[
 \int_{\mathbb{R}^N}\Big( |\Delta w_n(\tilde{x})|^2-\mu \frac{r_n^4}
 {\left| r_{n}\tilde{x}+y_n-a_j\right| ^4}w_n^2(\tilde{x})\Big)\,d\tilde{x}
 \geq \mathcal{A}_{\mu}^{N/4},
 \]
and so
 \[
 \int_{\mathbb{R}^N}\Big( |\Delta w_0(\tilde{x})|^2-\mu \frac{r_n^4}{|r_n
 \tilde{x}+y_n-a_j|^4}w_0^2(\tilde{x})\Big)\,d\tilde{x}
 \geq \mathcal{A}_{\mu}^{N/4}.
 \]
Consequently, 
 \begin{align*}
 J_{\mu}(u_n)
 &=\frac{1}{2}\int_{\Omega}(|\Delta u_n|^2-\mu V(x)u_n^2)\,dx
  -\frac{1}{2^*} \int_{\Omega} |u_n|^{2^*}\,dx+o(1)\\
 &=\frac{2}{N}\int_{\mathbb{R}^N}\left( |\Delta w_n(\tilde{x})|^2
  -\mu r_n^4 V(r_n\tilde{x}+y_n)w_n^2(\tilde{x})\right)\,d\tilde{x}+o(1)\\
 &\geq \frac{2}{N}\int_{\mathbb{R}^N}
 \Big( |\Delta w_n(\tilde{x})|^2-\mu r_n^4\frac{w_n^2(\tilde{x})}
 {|r_n \tilde{x}+y_n-a_j|^4}\Big)\,d\tilde{x}\\
 &\quad  +\frac{2}{N} \mu \int_{\mathbb{R}^N}r_n^4|r_n \tilde{x}+y_n-a_j
  |^{\alpha-4}w_n^2(\tilde{x})\,d\tilde{x}+o(1)\\
 &=\frac{2}{N}\int_{\mathbb{R}^N}\Big( |\Delta w_0(\tilde{x})|^2
  -\mu r_n^4\frac{w_0^2(\tilde{x})}{|r_n \tilde{x}+y_n-a_j|^4}\Big)\,d\tilde{x} \\
 &\quad +\frac{2}{N}\mu \int_{\mathbb{R}^N}\bar{r}^2|\bar{r}\tilde{x}+\bar{y}
 -a_j|^{\alpha-4}w_0^2(\tilde{x})\,d\tilde{x}\\
 &> \frac{2}{N}\mathcal{A}_{\mu}^{N/4},
 \end{align*}
which contradicts the assumption that 
\(J_{\mu}(u_n)\leq \frac{2}{N}\mathcal{A}_{\mu}^{N/4}\).
\smallskip

\noindent\textbf{Sub-step 2-2.}
 If \(\|\tilde{w}_n\|_{ \mathcal{D}^{2,2}(\mathbb{R}^N)}\to L>0\), set
 \[
 A+o(1)=\int_{\mathbb{R}^N}\left( |\Delta w_0(\tilde{x})|^2
 -\mu r_n^4V(r_n\tilde{x}+y_n)w_0^2(\tilde{x})\right) \, d\tilde{x}-\int _{\mathbb{R}^N}|w_0(\tilde{x})|^{2^*}\,d\tilde{x},
 \]
 then, by \eqref{d4} we obtain
\[
 -A+o(1)=\int_{\mathbb{R}^N}\left( |\Delta w_n(\tilde{x})|^2-\mu r_n^4V(r_n\tilde{x}+y_n)w_n^2(\tilde{x})\right) \, d\tilde{x}-\int _{\mathbb{R}^N}|w_n(\tilde{x})|^{2^*}\,d\tilde{x}.
\]
Suppose that \(A>0\) ($A<0$ can be considered similarly). We can find 
\(t_n\to 1, \) \(s_n\to 1\) such that
\( \bar {w}_n=t_n w_0\) and \(\bar {v}_n=s_n \tilde{w}_n\)
satisfy
 \begin{gather}\label{e1}
 A=\int_{\mathbb{R}^N}\left(|\Delta \bar {w}_n|^2-\mu r_n^4V(r_n\tilde{x}+y_n)
 \bar {w}_n^2\right) \,d\tilde{x}-\int _{\mathbb{R}^N}| \bar {w}_n|^{2^*}\,d\tilde{x},\\
\label{e2}
 -A= \int_{\mathbb{R}^N}\left(|\Delta \bar {v}_n|^2
 -\mu r_n^4V(r_n\tilde{x}+y_n)\bar {v}_n^2\right)\,d\tilde{x} 
 -\int _{\mathbb{R}^N}| \bar {v}_n|^{2^*}\,d\tilde{x}.
 \end{gather}
Now for \(v=t\bar{v}_n=t(s_n\tilde{w}_n), \ t\in (0,1)\), we have
 \[
 \int_{\mathbb{R}^N}\left(|\Delta v|^2-\mu r_n^4V(r_n\tilde{x}+y_n)v^2\right)\,d\tilde{x}=\int _{\mathbb{R}^N}|v|^{2^*}\,d\tilde{x}.
 \]
We denote
\[
 \tilde{J}_\mu(w)=\frac{1}{2}\int_{\mathbb{R}^N}\left(|\Delta w|^2-\mu r_n^4V(r_n\tilde{x}+y_n)w^2\right)\,d\tilde{x}-\frac{1}{2^*}\int _{\mathbb{R}^N}|w|^{2^*}\,d\tilde{x}.
\]
Then, we have
 \begin{align*}
 \tilde{J}_\mu(v)
 &=\frac{2}{N}\int_{\mathbb{R}^N}\left(|\Delta v|^2
  -\mu r_n^4V(r_n\tilde{x}+y_n)v^2\right)\,d\tilde{x}\\
 &=\frac{2}{N}t^2\int_{\mathbb{R}^N}\left(|\Delta \bar {v}_n|^2
  -\mu r_n^4V(r_n\tilde{x}+y_n)\bar {v}_n^2\right)\,d\tilde{x}\\
 &<\frac{2}{N}\int_{\mathbb{R}^N}\left(|\Delta \bar {v}_n|^2
  -\mu r_n^4V(r_n\tilde{x}+y_n)\bar {v}_n^2\right)\,d\tilde{x}\\
 &=\frac{1}{2}\int_{\mathbb{R}^N}\left(|\Delta \bar {v}_n|^2
  -\mu r_n^4V(r_n\tilde{x}+y_n)\bar {v}_n^2\right)\,d\tilde{x}
  -\frac{1}{2^*}\Big(\int _{\mathbb{R}^N}| \bar {v}_n|^{2^*}\,d\tilde{x}-A\Big) \\
 &=\tilde{J}_\mu(\bar {v}_n)+\frac{1}{2^*}A\\
 &=\tilde{J}_\mu(\bar {v}_n)+\frac{1}{2}A
  +\Big( \frac{1}{2^*}-\frac{1}{2}\Big) A\\
 &<\tilde{J}_\mu(w_0)+\tilde{J}_\mu( \tilde{w}_n)
  +\Big( \frac{1}{2^*}-\frac{1}{2}\Big) A+o(1).
 \end{align*}
Thus, we obtain
\[
 \tilde{J}_\mu(w_0)+\tilde{J}_\mu( \tilde{w}_n)>\tilde{J}_\mu(v)
 + \Big( \frac{1}{2}-\frac{1}{2^*}\Big) A.
 \]
Using an argument as in sub-step 2-1, we obtain
\[
 \tilde{J}_\mu(v)\geq \frac{2}{N}\mathcal{A}_{\mu}^{N/4}.
\]
It follows that
\begin{equation}\label{d6}
 \tilde{J}_\mu(w_0)+\tilde{J}_\mu( \tilde{w}_n)+o(1)
 \geq \frac{2}{N}\mathcal{A}_{\mu}^{N/4}
 + \Big( \frac{1}{2}-\frac{1}{2^*}\Big) A> \frac{2}{N}\mathcal{A}_{\mu}^{N/4}.
\end{equation}
On the other hand, we have
 \begin{align*}
 J_\mu(u_n)
 &=\frac{1}{2}\int_{\mathbb{R}^N}\left(|\Delta w_0|^2-\mu r_n^4V(r_n\tilde{x}+y_n)w_0^2\right)\,d\tilde{x}-\frac{1}{2^*}\int _{\mathbb{R}^N}|w_0|^{2^*}\,d\tilde{x}+o(1)\\
 &\quad +\frac{1}{2}\int_{\mathbb{R}^N}\left(|\Delta \tilde{w}_n|^2-\mu r_n^4V(r_n\tilde{x}+y_n) \tilde{w}_n^2\right)\,d\tilde{x}-\frac{1}{2^*}\int _{\mathbb{R}^N}| \tilde{w}_n|^{2^*}\,d\tilde{x}\\
 &=\tilde{J}_\mu(w_0)+\tilde{J}_\mu( \tilde{w}_n)+o(1).
 \end{align*}
 Using \eqref{d6}, we obtain
 \[
 J_\mu(u_n)> \frac{2}{N}\mathcal{A}_{\mu}^{N/4}.
 \]
 That is absurd in contrast to \eqref{d1}. When $A=0$, the proof is similar. 
 Using \eqref{e1}, \eqref{e2} with $A=0$, we obtain from the same argument as 
 above that
 \[
 J_\mu(u_n)=\tilde{J}_\mu(w_0)+\tilde{J}_\mu( \tilde{w}_n)+o(1)\geq \frac{2}{N}\mathcal{A}_{\mu}^{N/4}+\frac{2}{N}\mathcal{A}_{\mu}^{N/4}>\frac{2}{N}\mathcal{A}_{\mu}^{N/4}.
 \]
Again, we obtain a contradiction. 
\smallskip

\noindent\textbf{Step 3.} If \(r_n\to \bar{r}=0\) and \(\bar {y}\neq a_j\), then
\[
 r_n^4V(r_n\tilde{x}+y_n)\leq \frac{r_n^4}{\left|r_n\tilde{x}
 +y_n-a_j\right| ^4 }-r_n^4\left|r_n\tilde{x}+y_n-a_j\right|^{\alpha-4}\to 0.
 \]
Thus, we obtain
\begin{equation}\label{b1}
 J_\mu(u_n)=\frac{1}{2}\int_{\mathbb{R}^N}|\Delta {w}_n|^2d\tilde{x}
 -\frac{1}{2^*}\int _{\mathbb{R}^N}|{w}_n|^{2^*}d\tilde{x}+o(1).
\end{equation}
We also divide this step into two sub-steps.
\smallskip

\noindent\textbf{Sub-step 3.1.} 
When \(\|\tilde{w}_n\|_{ \mathcal{D}^{2,2}(\mathbb{R}^N)}\to 0\), we have
\[
 \int_{\mathbb{R}^N}|\Delta {w}_n|^2\,d\tilde{x}
 \geq \mathcal{A}_{0}\Big( \int _{\mathbb{R}^N}|{w}_n|^{2^*}\,d\tilde{x}\Big) ^{2/2^*}
 =\mathcal{A}_{0}\Big( \int_{\mathbb{R}^N}|\Delta {w}_n|^2\,d\tilde{x}\Big) ^{2/2^*},
\]
hence, we obtain
\[
 \int_{\mathbb{R}^N}|\Delta {w}_n|^2\,d\tilde{x}\geq \mathcal{A}_{0}^{N/4}.
\]
It follows from \eqref{b1} that
\begin{align*}
 J_\mu(u_n)
 &=\frac{1}{2}\int_{\mathbb{R}^N}|\Delta {w}_n|^2\,d\tilde{x}-\frac{1}{2^*}
 \int _{\mathbb{R}^N}|{w}_n|^{2^*}\,d\tilde{x}+o(1)\\
 &=\frac{2}{N}\int_{\mathbb{R}^N}|\Delta {w}_n|^2\,d\tilde{x}+o(1)\geq \frac{2}{N}\mathcal{A}_{0}^{N/4}> \frac{2}{N}\mathcal{A}_{\mu}^{N/4},
\end{align*}
which also contradicts \eqref{d1}.
\smallskip

\noindent\textbf{Sub-step 3.2.} 
If \(\|\tilde{w}_n\|_{ \mathcal{D}^{2,2}(\mathbb{R}^N)}\to L>0\), then the proof 
is similar. Using \eqref{d6}, we obtain from the same arguments as above that
 \begin{align*}
 J_\mu(u_n)
 &=\tilde{J}_\mu(w_0)+\tilde{J}_\mu( \tilde{w}_n)+o(1) \\
 &\geq  \begin{cases}
 \frac{2}{N}\mathcal{A}_{0}^{N/4}+ \left( \frac{1}{2}-\frac{1}{2^*}\right) A, 
  &\text{if } A\neq 0, \\
 \frac{2}{N}\mathcal{A}_{0}^{N/4}+\frac{2}{N}\mathcal{A}_{0}^{N/4},  
 &\text{if } A=0.
 \end{cases}
 \\
 &>\frac{2}{N}\mathcal{A}_{\mu}^{N/4}.
 \end{align*}
 Again, we obtain a contradiction.
 
 This completes the proof of Proposition \ref{p3}.
 \end{proof}
 
 \section{Proof of Theorem \ref{tm}}
 To obtain the existence of multiple solutions for the problem \eqref{11},
we need several lemmas.

 \begin{lemma}
For each \(u\in \mathcal{N}^-_j\) there exist a number \(\rho_u>0\) and a 
continuous function \(h>0\) defined on \(\{w \in H^2_0(\Omega): \|w\|<\rho_u \}\), 
satisfying 
\[
 h(0)=1, \ \ h(w)(u-w)\in \mathcal{N}^-_j , \ for \ \|w\|<\rho_u,
\]
and
\[
 \langle h'(0),w\rangle=\frac{2\int_{\Omega} \left( \Delta u\Delta w-\mu V(x)uw\right)\,dx -2^*\int_{\Omega} |u|^{2^*}uw\,dx-\lambda \int_{\Omega} fu\,dx}{\int_{\Omega} \left( |\Delta u|^2-\mu V(x)u^2\right)\,dx -(2^*-1)\int_{\Omega} |u|^{2^*}\,dx}.
\]
\end{lemma}

\begin{proof}
Define \(H:\mathbb{R} \times H^2_0(\Omega) \to \mathbb{R}\) as follows
\begin{align*}
 H(t,w)
 &=t\Big( \int_{\Omega}\left( |\Delta (u-w)|^2-\mu V(x)(u-w)^2\right)\,dx\Big)\\
 & \quad-t^{2^*-1}\int_{\Omega} |u-w|^{2^*}dx-\lambda \int_{\Omega} f(u-w)\,dx.
 \end{align*}
Since \(u\in \mathcal{N}^-_j\), we have
 \[
 H(1,0)= \int_{\Omega} \left( |\Delta u|^2-\mu V(x)u^2\right)\,dx -\int_{\Omega} |u|^{2^*}\,dx-\lambda \int_{\Omega} f u\,dx=0.
 \]
 and
 \[
 \frac{\partial H}{\partial t}(1,0)=\int_{\Omega} 
 \left( |\Delta u|^2-\mu V(x)(u)^2\right)\,dx-(2^*-1)\int_{\Omega} |u|^{2^*}\,dx \neq 0.
 \]
 We obtain the result by applying the implicit function Theorem at the point 
 \((1,0) \).
 \end{proof}
 
\begin{lemma}\label{l5}
Assume that {\rm (A1)--(A3)} are verified, and \(0<\mu<\bar{\mu}\). Then, for any 
sequence \(\{u_n^j\} \subset \mathcal{N}^-_j\), satisfying
 \[
 J_{\mu}(u_n^j)\to m<\frac{2}{N}\mathcal{A}_{\mu}^{N/4}-\chi \lambda^2,
 \]
 and
 \begin{equation}\label{p4}
 J'_{\mu}(u_n^j)\to 0, \ in \ H^{-2}(\Omega),
 \end{equation}
 \(\{u_n^j\}\) is relatively compact in \(H^{2}_0(\Omega)\).
 \end{lemma}
 
\begin{proof}
Since \(\{u_n^j\} \subset \mathcal{N}^-_j\) and \(J_{\mu}(u_n^j)\to m\) as 
\(n\to \infty\), we can assume from a similar argument to Proposition \ref{p3} that
\[
 0<\nu_3\leq \|u_n^j\|\leq \nu_4
\]
for some positive constants $\nu_3$ and $\nu_4$. Going if necessary to a subsequence, 
we can assume that \(u_n^j\rightharpoonup u^j\) in \(H^2_0(\Omega)\) and a.e.\
in \(\Omega\). Using \eqref{p4} then for any \(\varphi \in H^2_0(\Omega)\),
\begin{align*}
 o(1)&=\langle J'_{\mu}(u_n^j), \varphi \rangle\\
 &=\int _\Omega \left( \Delta u_n^j \Delta \varphi -\mu V(x)u_n^j \varphi \right)\,dx - \int _\Omega |u_n^j|^{2^*-2}u_n^j \varphi\, dx- \lambda \int _\Omega f\varphi \,dx\\
 &=\int _\Omega \left( \Delta u^j \Delta \varphi -\mu V(x)u^j \varphi \right)\,dx - \int _\Omega |u^j|^{2^*-2}u^j \varphi \,dx- \lambda \int _\Omega f\varphi \,dx +o(1)\\
 &=\langle J'_{\mu}(u^j), \varphi \rangle+o(1),
 \end{align*}
that is, \(u^j\) is a weak solution of \eqref{11}, and
\begin{equation}\label{z1}
 \int_{\Omega}\left( |\Delta u^j|^2-\mu V(x)(u^j)^2\right)\,dx - \int_{\Omega}|u^j|^{2^*}\,dx -\lambda \int_{\Omega}fu^j \,dx=0.
\end{equation}
Now, we claim that \(u^j\neq 0\). Arguing by contradiction, we assume that 
\(u^j\equiv 0\), then \(\lim_{n\to \infty}\int _\Omega fu_n^jdx
= \int _\Omega fu^jdx=0\) and therefore, from 
\(\{u_n^j\} \subset \mathcal{N}^-_j\) we have
\[
 \int_{\Omega}\left( |\Delta u_n^j|^2-\mu V(x)(u_n^j)^2\right)\,dx -\int_{\Omega} |u_n^j|^{2^*}\,dx=o(1),
\]
thus, we obtain
 \[
 \lim_{n\to \infty}\int_{\Omega}(|\Delta u_n^j|^2-\mu V(x)(u_n^j)^2)\,dx=\lim_{n\to \infty} \int_{\Omega} |u_n^j|^{2^*}\,dx=l_2>0,
 \]
and we find \(y_n\in \Omega, r_n\geq 0\) such that
\[
 \int _{B(y_n, r_n)}|u_n^j|^{2^*}\,dx=\frac{l_2}{2}.
\]
Again denoting \(w_n(\tilde{x})=r_n^{\frac{N-4}{2}}u_n^j(r_n\tilde{x}+y_n)\), by extending \(w_n(\tilde{x})\) to be zero for \(\tilde{x}\) outside \(\Omega_n\), we obtain \(w_n(\tilde{x})\in \mathcal{D}^{2,2}(\mathbb{R}^N)\) with
 \[
 \int_{\mathbb{R}^N}\left( |\Delta w_n(\tilde{x})|^2-r_n^4 V(r_n\tilde{x}+y_n)w_n^2(\tilde{x})\right) \,d\tilde{x}-\int _{\mathbb{R}^N}|w_n(\tilde{x})|^{2^*}\,d\tilde{x}=o(1).
 \]
 We divide the discussion into two cases.
\smallskip

\noindent\textbf{Case III.} 
If \(r_n\tilde{x}+y_n \notin B(a_j, \delta_0)\), it follows from the definition of 
\(\beta_j(u), \ \psi_{j}\) and \(r_0\) that
\[
 \frac{\delta_0}{3}=r_0=\beta_j(u_n^j)=\frac{\int_{\Omega} 
 \psi_{j}(x)|\Delta u_n^j(x)|^2\,dx}{\int_{\Omega} |\Delta u_n^j(x)|^2\,dx}
 =\frac{\int_{\mathbb{R}^N} \psi_{j}(r_n\tilde{x}+y_n)|\Delta w_n(\tilde{x})|^2
 \,d\tilde{x}}{\int_{\mathbb{R}^N} |\Delta w_n(\tilde{x})|^2\,d\tilde{x}} 
 \to \delta_0,
\]
which is a contradiction.
\smallskip

\noindent\textbf{Case IV.} If \(r_n\tilde{x}+y_n \in B(a_j, \delta_0)\), then
\begin{align*}
 J_{\mu}(u_n^j)&=\frac{1}{2}\int_{\Omega}(|\Delta u_n^j|^2-\mu V(x)(u_n^j)^2)\,dx-\frac{1}{2^*} \int_{\Omega} |u_n^j|^{2^*}\,dx+o(1)\\
 &=\frac{1}{2}\int_{\mathbb{R}^N}\left( |\Delta w_n(\tilde{x})|^2-\mu r_n^4 V(r_n\tilde{x}+y_n)w_n^2(\tilde{x})\right)\,d\tilde{x}-\frac{1}{2^*}\int _{\mathbb{R}^N}|w_n(\tilde{x})|^{2^*}\,d\tilde{x}+o(1)\\
 &=\frac{2}{N}\int_{\mathbb{R}^N}\left( |\Delta w_n(\tilde{x})|^2-\mu r_n^4 V(r_n\tilde{x}+y_n)w_n^2(\tilde{x})\right)\,d\tilde{x}+o(1)\\
 &=\tilde{J}_\mu(w_n)+o(1),
\end{align*}
The proof is divided into several steps:
step 1: \(r_n\to \bar{r}=0\) and \(\bar {y}=a_j\);
step 2: \(r_n\to \bar{r}>0\);
step 3: \(r_n\to \bar{r}=0\) and \(\bar {y}\neq a_j\).

We use the same arguments as those in the corresponding steps in the proof of 
Proposition \ref{p3} to get that 
$\tilde{J}_\mu(w_n)\geq \frac{2}{N}\mathcal{A}_{\mu}^{N/4}$.
Here, we only sketch the argument for step 2. 
It follows from \eqref{A}, that
\begin{align*}
 &\int_{\mathbb{R}^N}\Big( |\Delta w_n(\tilde{x})|^2-\mu \frac{r_n^4}{\left| 
  r_{n}\tilde{x}+y_n-a_j\right| ^4}w_n^2(\tilde{x})\Big)\,d\tilde{x}\\
 &\geq \mathcal{A}_{\mu} \Big( \int _{\mathbb{R}^N}|w_n(\tilde{x})|^{2^*}\,d\tilde{x}
  \Big) ^{2/2^*}\\
 &=\mathcal{A}_{\mu} \Big(\int_{\mathbb{R}^N} 
  \Big( |\Delta w_n(\tilde{x})|^2-\mu r_n^4 V(r_n\tilde{x}+y_n)w_n^2(\tilde{x})\Big) 
  \,d\tilde{x}\Big)^{2/2^*}\\
 &\geq \mathcal{A}_{\mu} \Big(\int_{\mathbb{R}^N}
 \Big( |\Delta w_n(\tilde{x})|^2-\mu \frac{r_n^4}{\left| r_{n}\tilde{x}
  +y_n-a_j\right| ^4}w_n^2(\tilde{x})\Big)\, d\tilde{x}\Big)^{2/2^*}. 
\end{align*}
So, that
\[
 \int_{\mathbb{R}^N}\Big( |\Delta w_n(\tilde{x})|^2-\mu \frac{r_n^4}
 {\left| r_{n}\tilde{x}+y_n-a_j\right| ^4}w_n^2(\tilde{x})\Big)\,d\tilde{x}
 \geq \mathcal{A}_{\mu}^{N/4}.
\]
Thus, from (A2) it follows that
\[
 \int_{\mathbb{R}^N}\left( |\Delta w_n(\tilde{x})|^2-\mu r_n^4 V(r_n\tilde{x}+y_n)w_n^2(\tilde{x})\right)\, d\tilde{x}
 \geq \mathcal{A}_{\mu}^{N/4}.
\]
Then, we have
 \[
 \frac{2}{N}\mathcal{A}_{\mu}^{N/4}-\chi \lambda^2> m
 =\lim_{n\to \infty}J_\mu (u_n^j)=\lim_{n\to \infty}\tilde{J}_\mu(w_n)
 \geq \frac{2}{N}\mathcal{A}_{\mu}^{N/4},
 \]
 which is a contradiction. In sum, we have proved that \(u^j\neq 0\).
 
Next, we prove that the limit \(u^j\) is indeed strong. Suppose by contradiction that 
\(\|u_n^j-u^j\|\to \nu_5> 0\) and denote \(v_n=u_n^j-u^j\), then we have 
\(v_n\rightharpoonup 0\) in \(H^2_0(\Omega)$ and $\|v_n\|\to \nu_5> 0\).
It follows from Brezis-Lieb Lemma and \eqref{z1}, that
 \begin{align*}
&\int_{\Omega}\left( |\Delta v_n(x)|^2- V(x)v_n^2(x)\right)\, dx
 -\int_{\Omega}|v_n(x)|^{2^*}\,dx\\
&=\int_{\Omega}\left( |\Delta u^j_n(x)|^2-V(x)(u^j_n)^2(x)\right) \,dx
 -\int _{\Omega}|u^j_n(x)|^{2^*}\,dx\\
& \quad - \int_{\Omega}\left( |\Delta u^j(x)|^2-r_n^4 V(x)(u^j)^2(x)\right) \,dx
 +\int _{\Omega}|u^j(x)|^{2^*}\,dx+o(1)\\
&=\lambda \int_{\Omega} fu_n^j\,dx-\lambda \int_{\Omega} fu^j\,dx+o(1)\\
&=o(1).
\end{align*}
Similarly to the previous proof, we can find \(y_n\in \Omega\), \(r_n>0\) such that 
\[
 \int _{B(y_n, r_n)}|u_n^j|^{2^*}\,dx=\frac{l_3}{2}>0.
 \]
Letting \(z_n(\tilde{x})=r_n^{\frac{N-4}{2}}v_n(r_n\tilde{x}+y_n)\), 
and extending \(z_n\) to be zero for \(\tilde{x}\) outside 
\(\Omega_n=r_n^{-1}(\Omega-y_n)\), we can write 
\(z_n(\tilde{x})\in D^{2,2}(\mathbb{R}^N)\) and we have that
\begin{equation}\label{z2}
 \int_{\mathbb{R}^N}\left( |\Delta z_n(\tilde{x})|^2-r_n^4 V(r_n\tilde{x}+y_n)z_n^2(\tilde{x})\right)\, d\tilde{x}-\int _{\mathbb{R}^N}|z_n(\tilde{x})|^{2^*}\,d\tilde{x}=o(1).
\end{equation}
We will distinguish two cases.
\smallskip

\noindent \textbf{Case V.} 
If \(r_n\tilde{x}+y_n \notin B(a_j, \delta_0)\), we can use similar arguments as 
those in case III to obtain a contradiction,
so we omit the details here.
\smallskip

\noindent\textbf{Case VI.} If \(r_n\tilde{x}+y_n \in B(a_j, \delta_0)\), then
 \begin{align*}
 J_{\mu}(u_n^j)
 &=\frac{1}{2}\int_{\Omega}(|\Delta v_n|^2-\mu V(x)v_n^2)\,dx-\frac{1}{2^*} \int_{\Omega} |v_n|^{2^*}\,dx+J_{\mu}(u^j)+o(1)\\
 &=\frac{1}{2}\int_{\mathbb{R}^N}\left( |\Delta z_n(\tilde{x})|^2
 -\mu r_n^4 V(r_n\tilde{x}+y_n)z_n^2(\tilde{x})\right)\,d\tilde{x}-\frac{1}{2^*}
 \int _{\mathbb{R}^N}|z_n(\tilde{x})|^{2^*}\,d\tilde{x}
  +J_{\mu}(u^j)+o(1)\\
 &=\tilde{J}_\mu(z_n)+J_{\mu}(u^j)+o(1).
 \end{align*}
 Now, we need to distinguish three steps.
\smallskip

\noindent\textbf{Step 1.} If \(r_n\to \bar{r}>0\),
it follows from \eqref{A} and \eqref{z2}, that
 \begin{align*}
&\int_{\mathbb{R}^N}\Big( |\Delta z_n(\tilde{x})|^2-\mu \frac{r_n^4}
{\left| r_{n}\tilde{x}+y_n-a_j\right| ^4}z_n^2(\tilde{x})\Big)\,d\tilde{x}\\
 &\geq \mathcal{A}_{\mu} \Big( \int _{\mathbb{R}^N}|z_n(\tilde{x})|^{2^*}\Big) ^{2/2^*}\,d\tilde{x}\\
 &\geq \mathcal{A}_{\mu} \Big(\int_{\mathbb{R}^N} |\Delta z_n(\tilde{x})|^2
 -\mu \frac{r_n^4}{\left| r_{n}\tilde{x}+y_n-a_j\right| ^4}z_n^2(\tilde{x})
 \Big)^{2/2^*} \,d\tilde{x}.
 \end{align*}
Thus, we obtain
\begin{equation}\label{z3}
 \tilde{J}_\mu(z_n)=\frac{2}{N}\int_{\mathbb{R}^N}
 \left( |\Delta z_n(\tilde{x})|^2-\mu r_n^4 V(r_n\tilde{x}+y_n)z_n^2(\tilde{x})
 \right)\,d\tilde{x}\geq \frac{2}{N} \mathcal{A}_{\mu}^{N/4}.
 \end{equation}
 
\noindent\textbf{Step 2.} If \(r_n\to \bar{r}=0\) and \(\bar {y}\neq a_j\) then we have
 \[
 r_n^4V(r_n\tilde{x}+y_n)\leq \frac{r_n^4}{\left|r_n\tilde{x}+y_n-a_j\right| ^4 }
 -r_n^4\left|r_n\tilde{x}+y_n-a_j\right|^{\alpha-4}\to 0.
 \]
Thus, we obtain
 \(\tilde{J}_\mu(z_n)=\frac{2}{N}\int_{\mathbb{R}^N} 
 |\Delta z_n(\tilde{x})|^2\,d\tilde{x}\). Then, we obtain
\[
 \int_{\mathbb{R}^N}|\Delta {z}_n|^2\,d\tilde{x}
 \geq \mathcal{A}_{0}\Big( \int _{\mathbb{R}^N}|{z}_n|^{2^*}\,d\tilde{x}\Big) ^{2/2^*}
 =\mathcal{A}_{0}\Big( \int_{\mathbb{R}^N}|\Delta {z}_n|^2\,d\tilde{x}\Big) ^{2/2^*},
\]
hence
 \begin{equation}\label{z4}
 \tilde{J}_\mu(z_n)\geq \frac{2}{N} \mathcal{A}_{0}^{N/4}\geq \frac{2}{N} \mathcal{A}_{\mu}^{N/4}.
 \end{equation}
 
\noindent\textbf{Step 3.}
 If \(r_n\to \bar{r}=0\) and \(\bar {y}=a_j\), then similarly, we can obtain
 \begin{equation}\label{z5}
 \tilde{J}_\mu(z_n)\geq \frac{2}{N} \mathcal{A}_{0}^{N/4}
 \geq \frac{2}{N} \mathcal{A}_{\mu}^{N/4}.
 \end{equation}
It follows from \eqref{z3}, \eqref{z4} and \eqref{z5} that
\[
 \tilde{J}_\mu(z_n)\geq \frac{2}{N} \mathcal{A}_{\mu}^{N/4}.
\]
Since $u^j$ is a solution of \eqref{11} then from Lemma \ref{lm}, we obtain
 \[
 J_{\mu}(u^j)\geq -\chi \lambda^2.
 \]
 Therefore, 
 \[
 J_\mu(u^j_n)=\tilde{J}_\mu(z_n)+J_{\mu}(u^j)+o(1)
 \geq \frac{2}{N} \mathcal{A}_{\mu}^{N/4} -\chi \lambda^2,
 \]
which is a contradiction with the choice of \(\{u_n^j\}\).

This completes the proof of Lemma \ref{l5}.
\end{proof}

\begin{proof}[Proof of Theorem \ref{tm}]
 Fix $\lambda^*=\min\{\lambda_1, \lambda_2, \lambda_3\}$.
For any $j\in \{1,\dots,k\}$, we obtain 
$\overline{\mathcal{N}_j^-}=\mathcal{N}_j^-\cup \Gamma_j^-$. 
Propositions \ref{PP} and \ref{p3} imply that
$$
 c_{\lambda, j}<\frac{2}{N}\mathcal{A}_{\mu}^{N/4}-\chi \lambda^2< \hat{c}_{\lambda, j}.
$$
Thus, $c_{\lambda, j}=\inf \{J_{\mu}(u), u\in \overline{\mathcal{N}_j^-}\}$. 
By the Ekeland variational principle, we can obtain the minimizing sequence 
$\{u_n^j\}\subset \overline{\mathcal{N}_j^-}$ satisfying
 $$
 c_{\lambda, j} <J_{\mu}(u_n^j)<c_{\lambda, j}+\frac{1}{n}\quad\text{and}\quad
 J'_{\mu}(u_n^j)\to 0, \ in \ H^{-2}(\Omega).
 $$
Therefore, by Lemma \ref{l5} up to a subsequence we have $u_n^j\to u^j$ and $u^j\neq 0$ 
in $H^2_0(\Omega)$ and so $u^j$ is a solution of \eqref{11}.
Moreover, from Remark \ref{RM}, we know that $u^i$ and $u^j$ are distinct if 
$i\neq j.$ This implies that problem \eqref{11} has at least $k$ solutions 
$u^j\in \mathcal{N}_j^-$.
\end{proof}

\begin{thebibliography}{10}

\bibitem{c1} J. Q. Chen;
\emph{Multiple positive solutions for a semilinear equation with prescribed singularity}, 
 J. Math. Anal. Appl. Vol. 305, pp. 140-157, 2005.

\bibitem{j1} L. D'Ambrosio, E. Jannelli;
\emph{Nonlinear critical problems for the biharmonic operator with Hardy potential},
 Calc. Var. Partial Differ. Equ. Vol. 54, pp. 365-396, 2015.

\bibitem{ek} K. L. Ekinci, M. L. Roukes;
\emph{Nanoelectromechanical system}s, Rev. Sci. Instrum. Vol. 76, 2005.

\bibitem{g} F. Gazzola, H. C. Grunau, G. Sweers;
\emph{Polyharmonic Boundary Value Problems: Positivity Preserving and Nonlinear Higher 
Order Elliptic Equations in Bounded Domains}, Springer, 2010.

\bibitem{gw} S. Golubovi\'c, A. Wang;
\emph{Capillary surface growth and thin‑film instabilities modeled by higher-order 
PDEs}, Phys. Rev. Lett. Vol. 69, pp. 2535-2538, 1992.

\bibitem{hw} Y. Huang;
\emph{Multiple positive solutions of nonhomogeneous equation involving the p-Laplacian}, 
Nonlinear Anal. Vol. 43, pp. 905-922, 2001.

\bibitem{b5} D. Kang, P. Xiong;
\emph{Ground state solutions to biharmonic equations involving critical nonlinearities 
and multiple singular potentials}, Appl. Math. Lett. Vol. 66, pp. 9-15, 2017.

\bibitem{K} D. Kang, L. Xu;
\emph{Asymptotic behavior and existence results for biharmonic problems involving 
Rellich potentials}, J. Math. Anal. Appl. Vol. 455, pp. 1365-1382, 2017.

\bibitem{lin} K. Lindsay, J. Quaife, E. Wendelberger;
\emph{Numerical study of thin plate vibration with multiple clamped points: 
singular eigenmodes and mode elimination}, J. Vib. Acoust. Vol. 139, 2017.

\bibitem{bh} J. L. F. Melo, E. M. Santos;
\emph{Positive solutions to a fourth-order elliptic problem by the Lusternik-Schnirelmann 
category}, J. Math. Anal. Appl. Vol. 420, pp. 532-550, 2014.

\bibitem{bt} G. Li, T. Yang, L. Huang;
\emph{Existence of two weak solutions for biharmonic equations with Hardy-Sobolev 
critical exponent and non-homogeneous perturbation term} (in Chinese),
 Sci. Sin. Math. Vol. 49, pp. 1813-844, 2019.

\bibitem{met} Y. Mao, R. Huang, C. Yang;
\emph{Elastic multipole model for holes/inclusions in elastic media: stress 
interactions via multipolar fields in metamaterials}, Phys. Rev. E Vol. 103, 2021.

\bibitem{ms} S. Messirdi, A. Bennour, A. Matallah;
\emph{On nonhomogeneous biharmonic problems involving Rellich-type potentials and 
a critical Sobolev exponent}, Indian J. Pure Appl. Math. Vol. 54, pp. 110-124, 2023.

\bibitem{bh2} X. Qian, J. Wang, M. Zhu;
\emph{Multiple nontrivial solutions for a class of biharmonic elliptic equations 
with Sobolev critical exponent}, Math. Probl. Eng, 2018.

\bibitem{t} G. Tarantello;
\emph{On nonhomogeneous elliptic equations involving critical Sobolev exponent}, 
Ann. Inst. H. Poincar\'e Anal. Non Lin\'eaire. Vol. 9, pp. 281-304, 1992.

\bibitem{hz} J. Zhang, T. S. Hsu;
\emph{Multiplicity results for biharmonic equations involving multiple Rellich-type 
potentials and critical exponents}, Bound. Value Probl. Vol. 2019, pp. 103-122, 2019.

\end{thebibliography}


\end{document}
